\textbf{Grokking.} Power et al.~\cite{power2022grokking} introduced grokking on small algorithmic datasets including modular arithmetic and reported that data fraction and weight decay govern the delay. Nanda et al.~\cite{nanda2023progress} reverse-engineered the one-layer transformer solution for modular addition and defined progress measures that change smoothly before the sudden test-accuracy jump. Liu et al.~\cite{liu2022understanding} give an effective theory and phase diagrams over hyper-parameters, and Liu et al.~\cite{liu2022omnigrok} relate grokking to a mismatch between the training and test loss landscapes as a function of weight norm. Gromov~\cite{gromov2023grokking} shows grokking in two-layer MLPs on modular tasks, and Mohamadi et al.~\cite{mohamadi2024why} prove that in the kernel regime a constant fraction of all pairs is needed for permutation-equivariant models. Golechha~\cite{golechha2024progress} studies grokking on real-world tasks and questions the weight-norm explanation. None of these studies, to our knowledge, manipulates the \emph{order} in which training examples are presented, which is the gap this study targets (we do not claim to have exhaustively searched the literature).

\textbf{Curriculum learning.} Bengio et al.~\cite{bengio2009curriculum} proposed presenting examples from easy to hard; the survey of Wang et al.~\cite{wang2020survey} organises later work into difficulty measurers and training schedulers. Our curriculum follows that template with a fixed difficulty measure (operand size) and a linear pacing function, and the anti-curriculum and shuffled-order controls are the usual ablations of the difficulty measurer.

\textbf{Optimiser and architecture.} We use AdamW with decoupled weight decay~\cite{loshchilov2017decoupled} and a standard transformer block~\cite{vaswani2017attention}.
